% sections/ndss_07_discussion.tex  — NDSS 2027

\subsection{Limitations}

\paragraph{LSCPS as lower bound.}
Our rule-based capability extraction (\lscps) is a conservative
precision-recall tradeoff: it minimizes false positives at the cost
of missing obfuscated or dynamically-constructed dangerous calls.
Reported IC rates are therefore \emph{lower bounds} on true inconsistency.
CR = \MeanCR is an \emph{upper bound} on true disclosure rates.
LLM-based semantic adjudication (deferred to future work)
would expand detection coverage.

\paragraph{CAPE framework limitations.}
Of 85 pairs attempted at submission, 24 timed out under our pre-registered
175\,s hard-kill policy (LLM context compaction at high token usage).
We report both ITT (67.1\%) and completed-run (93.4\%) metrics side-by-side.
Controls and P2/P3 payload validations are 20-pair and 8-pair spot-checks;
expanding these to the full 125 pairs would further strengthen the causal claim.
We have not exhaustively tested all possible injection phrasings;
adversarial rephrasing likely increases the success rate further.

\paragraph{Single platform.}
Our measurement covers \clawhub and skills.sh only.
Other agent platforms (Cursor, Continue.dev, Claude Computer Use)
may have different transparency practices.
Cross-platform validation of skill.sh (90.9\% inconsistency)
suggests the finding generalizes.

\paragraph{Temporal snapshot.}
The ClawHub dataset reflects March 2026.
Skill descriptions may improve over time, though our temporal
analysis (\S\ref{sec:temporal}) shows the trend is toward
capability accretion without description update.

\subsection{Defenses}

\paragraph{Capability manifests.}
Marketplace platforms should require machine-readable capability
declarations (analogous to Android permissions) listing:
file-system access scope, network destinations, shell invocation flags,
credential access, and external API calls.
Our \lscps taxonomy provides a starting point for this schema.

\paragraph{Publish-time screening.}
Our \sysname pipeline can run at publish time to flag skills with
CR below a threshold (\eg CR $<$ 0.5), preventing the most egregious
mismatch patterns from reaching users.
At our measured FPR of $<$2\%, the cost is acceptable for platform operators.

\paragraph{Runtime capability monitoring.}
Agent runtimes should log and optionally confirm dangerous tool calls
(\texttt{exec}, \texttt{write}, \texttt{delete}) even when initiated by
installed skills, with user-visible notifications.
Our results show the agent's current safety model is heuristic-based
and inconsistent: shell exec is fully unguarded while obvious memory-write
requests are refused. A consistent capability-based policy is needed.

\paragraph{Trust-building attack countermeasure.}
The session warm-up bypass (\S\ref{sec:cape}) suggests that
agent safety policies should not rely on conversational trust context.
Capability whitelists established at skill-install time---not inferred
from session history---would prevent this bypass.

\subsection{Broader Implications}

The IC-1 finding (\ICone of dangerous skills omit capabilities) implies
that LLM agent marketplaces are currently operating on an \emph{honor system}
for capability disclosure.
As agent capabilities grow---computer use, multi-agent orchestration,
financial autonomy---the consequences of this trust gap compound.
Our work provides both the measurement evidence and the attack validation
needed to motivate platform-level policy change.
